\section{Dyskusja wyników}
Przeprowadzone w niniejszej pracy badania pozwoliły na wieloaspektową analizę metod mitygacji podatności dużych modeli językowych na bezpośrednie ataki promptowe. Dokonane pomiary użyteczności modelu oraz opóźnień wprowadzanych przez techniki obrony umożliwiły osadzenie problematyki w kontekście szerszym niż wyłącznie minimalizacja skuteczności ataków. Zgromadzone rezultaty wskazują, że dążenie do redukcji \gls{ASR} może być związane ze spadkiem interaktywności systemu lub negatywnym wpływem na obsługę standardowych zapytań.

Zjawisko to jest szczególnie dostrzegalne na bazie wyników otrzymanych dla metody wykorzystującej zewnętrzny klasyfikator Llama-Guard. Zapewniał on zdecydowanie najlepszą ochronę przed rozpatrywanymi atakami. Dla każdego z nich otrzymano najniższy poziom wskaźnika \gls{ASR}, niezależnie od rodziny modelu. Należy jednak dostrzec narzut tej metody na system. Wykorzystanie modelu Llama-Guard-3-8B, nawet przy zastosowaniu kwantyzacji do postaci 8-bitowej, było związane z koniecznością alokacji dodatkowych 8 GB pamięci VRAM. Może stanowić to przeszkodę dla środowisk z ograniczonymi zasobami. Natomiast inne badane metody, takie jak Self-Reminder oraz sterowanie aktywacjami, nie wymagały wczytywania dodatkowych modeli oraz nie wymagały tak znaczącego narzutu pamięciowego. Pomimo najwyższej skuteczności obrony, klasyfikator Llama-Guard wywołał również najwięcej nadmiernych odmów na bezpieczne polecenia. Z tego samego względu dla tej metody zaobserwowano spadek poprawności udzielonych odpowiedzi w teście sprawdzającym ogólne zdolności modeli. Co więcej, dla metody opartej na Llama-Guard wyłącznie na wejściu otrzymano ponad dziesięciokrotnie wyższy czas wygenerowania pierwszego tokena od pozostałych metod (ponad 400 ms). Może to negatywnie wpływać na doświadczenie użytkownika systemu. Wdrożenie ochrony modelu docelowego na wejściu i wyjściu przy użyciu Llama-Guard uniemożliwiło strumieniowanie. Czas oczekiwania na zwrócenie odpowiedzi wynosił około 5 s, co zdecydowanie degraduje interaktywność systemu.

Znacznie lepszą alternatywą pod kątem minimalizacji wprowadzanych opóźnień i wykorzystania zasobów okazała się metoda obrony oparta o sterowanie aktywacjami. Wprowadziła znikomy narzut na czas otrzymania pierwszego tokenu rzędu milisekundy względem poziomu bazowego, a ponadto zachowała bazową przepustowość generacji kolejnych tokenów. Metoda ta osiągnęła również mniejszy o 5 punktów procentowych poziom wskaźnika \gls{ORR} niż zewnętrzny klasyfikator w przypadku modelu Llama-3.1-8B-Instruct oraz aż o 22 punkty procentowe mniej dla modelu Mistral-7B-Instruct-v0.3. Ponadto, w przypadku modelu z rodziny Llama, dla metody mitygacji opartej o sterowanie aktywacjami uzyskano zbliżoną do Llama-Guard skuteczność obrony. Osiągnięcie to zostało dokonane przy jednoczesnym zachowaniu obniżonego poziomu nadmiernej odmowy, brakiem wpływu na rezultat w teście MMLU oraz zachowaniem wysokiego poziomu responsywności systemu.

Niski stopień opóźnień, porównywalny ze sterowaniem aktywacjami, uzyskano przy użyciu metody obrony Self-Reminder. Dostrzeżono również dla tego podejścia nieznaczny wzrost wskaźnika \gls{ORR} względem poziomu bazowego. Metoda ta natomiast okazała się niewystarczająca pod względem ochrony modelu Llama, ponieważ w aż 43\% przypadków uległa atakowi \gls{DAN} oraz uzyskała najgorszy wynik przeciwko metodzie \gls{PAIR}. Co więcej, analiza wyników testu MMLU dowiodła, że metoda Self-Reminder w największym stopniu wpłynęła negatywnie na liczbę poprawnie wybranych odpowiedzi w porównaniu do innych metod. Przetwarzania instrukcji bezpieczeństwa w oknie kontekstowym mogło w niektórych przypadkach spowodować degradację zdolności precyzyjnego wnioskowania modelu. Osiągnięte dla techniki Self-Reminder wyniki wskazują, że inżynieria promptów stosowana samodzielnie nie zapewnia odpowiedniego poziomu bezpieczeństwa.

Należy zauważyć, że zadowalającą redukcję skuteczności ataków przy zastosowaniu sterowania aktywacjami zaobserwowano wyłącznie dla modelu Llama-3.1-8B-Instruct. Skuteczność tej metody jest zależna od podstawowego bezpieczeństwa modelu, czyli od metod użytych na etapie jego trenowania. Na podstawie porównania otrzymanych rezultatów dla modelu Llama oraz Mistral można stwierdzić, że sterowanie aktywacjami wymaga, by model posiadał wysoki stopień wbudowanego bezpieczeństwa. Nie jest to zatem technika uniwersalna, którą można stosować z pozytywnym efektem w ramach dowolnego modelu. Wektor sterujący obliczany w ramach tej metody musi być wyliczony dla każdego poszczególnego modelu. Wymaga to przeprowadzenia procesu ekstrakcji osobno dla każdej architektury. Skuteczność tej metody zależna jest również w dużej mierze od doboru hiperparametrów, które również należy dla każdego modelu dopasować indywidualnie. Wpływ na działanie metody miał wybór warstwy, na której znajdowała się bramka sterująca, a także współczynnika $\alpha$ określającego siłę interwencji. Ich nieoptymalny wybór może negatywnie oddziaływać na finalne bezpieczeństwo systemu. Podczas eksperymentów zauważono, że zbyt duża wartość $\alpha$ może sprawić, że model generuje zaburzone odpowiedzi. Funkcjonowanie tej metody mitygacji może być również uzależnione od doboru zbioru danych par promptów o przeciwstawnej intencji. Pod tym względem zewnętrzny klasyfikator zachowuje przewagę, ponieważ zapewnia zbliżony, wysoki poziom bezpieczeństwa niezależnie od rodziny chronionego modelu.